\documentclass[10pt]{article}

\usepackage[utf8]{inputenc}

\usepackage[T1]{fontenc}

\usepackage{amsmath}

\usepackage{amsfonts}

\usepackage{amssymb}

\usepackage[version=4]{mhchem}

\usepackage{stmaryrd}

\usepackage{bbold}

\usepackage{mathrsfs}



\begin{document}

Лum.: [1] Кля цк и н В. И., Т а т а рск и й В. И., «Успехи физич. наук», 1973, т. 110, № 4, с. 499-536; [2] Монин А. С., Я гл о м А. М., Статистическая гидромеханика. Механика турбулентности, ч. 2, М., 1967; [3] Edwards S. F., «J. Fluid Mech.», 1964, v. 18, № 2, p. 239-73; [4] He rri ng J. H., "Phys. Fluids», 1965, v. 8, № 12, p. 2219-25; [5] Kraichnan R. H., «J. Fluid Mech.», 1971, v. 47, № 3, p. 513-24. A. П. Мирабель.



ЛАПЛАСА ДВУСТОРОННЕЕ ПРЕОБРАЗОВАНИЕ см. Фурье интеграл.



ЛАПЛАСА ИНТЕГРАЛ - см. Лапласа преобразование.



ЛАПЛАСА OПEPATOP, л а П л а с и а н,- простейший эллиптический дифференциальный оператор 2-го порядка:



\[

\Delta f=\frac{\partial^{2} f}{\partial x_{1}^{2}}+\ldots+\frac{\partial^{2} f}{\partial x_{n}^{2}},

\]



действуюший на гладкие функции $f\left(x_{1}, \ldots, x_{n}\right)$, определенные в евклидовом пространстве $\mathbb{R}^{n}$ с декартовыми координатами $x_{1}, \ldots, x_{n}$ (или в некоторой его части $G$ ). Л. о. инвариантен относительно ортогональных преобразований координат в $\mathbb{R}^{n}$, то есть преобразований



\[

x_{i}^{\prime}=\sum_{k} s_{i k} x_{k}

\]



с ортогональной матрицей $s_{i k}$.



ЛАПЛАСА ПРЕОБРАЗОВАНИЕ - преобразование, переводящее функщию $f(x)$ действительного переменного $t$, $0<t<\infty$, называемую о р и г и на л о м, в функцию



\[

F(p)=L[f]=\int_{0}^{\infty} f(t) e^{-p t} d t

\]



комплексного переменного $p=\sigma+i \tau$. Под Л. П. понимают также не только само преобразование, но и его результат функцию $F(p)$. Интеграл в правой части формулы (1) называется интегралом Лапласа. Он был рассмотрен П. Лапласом (P. Laplace) в ряде работ, к-рые объединены в его книге «Аналитическая теория вероятностей», вышедшей в 1812. Значительно раныше (1737) такие интегралы применял к решению дифференциальных уравнений Л. Эйлер (L. Euler).



При нек-рых условиях Л. п. определяют функцию $f(t)$ однозначно, в простейших случаях - по формуле обр а ш е и я:



\[

f(t)=\frac{1}{2 \pi i} \int_{\gamma-i \infty}^{\gamma+i \infty} F(p) e^{p t} d t

\]



Л. п. является линейным функциональным преобразованием. Из числа основных формул Л. П. можно отметить следующие:



\[

\begin{gathered}

L\left[f^{\prime}(t)\right]=p F(p)-f(0), \\

L\left[t^{n} f(t)\right]=(-1)^{n} F^{(n)}(p), \quad n=1,2, \ldots, \\

L\left[\int_{0}^{t} f(t) d t\right]=\frac{F(p)}{p}, \quad t>0 .

\end{gathered}

\]



Л. п. в сочетании с формулой (2) его обращения применяется к интегрированию дифференциальных уравнений. В частности, в силу свойства (1) и линейности Л. П. решение обыкновенного линейного дифференциального уравнения с постоянными коэффициентами удовлетворяет алгебраич. уравнению 1-й степени и может быть, следовательно, легко найдено. Так, если, напр., $y^{\prime \prime}+y=f(t), y(0)=y^{\prime}(0)=0$ и $Y(p)=L[y], F(p)=L[f]$, то $L\left[y^{\prime \prime}\right]=p^{2} Y(p)$ и $p^{2} Y(p)+$ $+Y(p)=F(p)$, откуда $Y(p)=\frac{F(p)}{p^{2}+1}$.



Многочисленные задачи электротехники, гидродинамики, механики, теплопроводности эффективно решаются методами, использующими Л. п.



Лum.: [1] Д е ч Г., Руководство к практическому применению преобразований Лапласа, пер. с нем., М., 1965. ЛАПЛАСА ПPEOБPAЗОВАНИЕ о б о бщ е н о й Ф у н к ц и и - преобразование Лапласа, определенное для обобценных функций, основное пространство которых содержит ли н е й ны е экспон енты - функции вида



\[

e^{i z}(\zeta)=e^{i z \zeta}

\]



$\zeta$ - вещественный или комплексный аргумент, $z$ - параметр, лежаший в трубчатой области $T^{O}=\mathbb{R}^{n}+i O \subset \mathbb{C}^{n}, O-$ область в $\mathbb{R}^{n}, z \zeta=z_{1} \zeta_{1}+\ldots+z_{n} \zeta_{n}$. В этом случае Л. П. $L[g]$ обобшенной функции $g$ определяется формулой



\[

L[g](z)=\left(g, e^{i z}\right), z \in T^{O} .

\]



В нормальных условиях топология основного пространства обеспечивает аналитичность функции $f(z)=L[g](z)$ в $T^{O}$.



Л. П. обобщенных функций было введено Л. Шварцем (L. Schwartz) и Ж. Лионсом (J. Lions) в начале 50-х гг. 20 в. Оно широко применяется в аксиоматич. квантовой теории поля, в теории дифференциальных уравнений с частными производными, в тауберовой теории для обобщенных функций, при изучении граничных значений аналитич. Функций и в других областях математики и математич. физики.



Л. П. аналитических функционалов. Основное пространство $x(O), O-$ область в $\mathbb{C}^{n}$, содержит все линейные экспоненты, так что для всякого аналитич. функционала $g \in x^{\prime}(O)$ определено Л. П. $L[g](z), z \in \mathbb{C}^{n},-$ целая функция экспоненциального типа. Подробнее, пусть



\[

E(O)=\operatorname{limind}_{K \subset O} E(K),

\]



где $K$ - комплексный компакт, $E(K)$ - банахово пространство всех целых функций $f(z)$ с конечной нормой



\[

|f|=\sup \left\{|f(z)| e^{-H_{K}(z)} ; z \in \mathbb{C}^{n}\right\},

\]



где $H_{K}(z)=\sup \{\operatorname{Re} i z \zeta ; \zeta \in K\}-$ опорная функция $K$. Очевидно, $E(O)=E(\operatorname{ch} O)$, где $\operatorname{ch} O-$ выпуклая оболочка $O$. Л. П. $L$ есть линейное непрерывное отображение $x^{\prime}(O)$ в $E(O)$. Если область $O$ выпуклая, то



\[

L: x^{\prime}(O) \cong E(O)

\]



\begin{itemize}

  \item топологич. изоморфизм.

\end{itemize}



Л. П. вешественных аналитических функцио налов. Основное пространство $\mathscr{A}(\omega), \omega-$ область в $\mathbb{R}^{n}$, содержит все линейные экспоненты, то есть и здесь для всякого $g \in \mathscr{A}^{\prime}(\omega)$ определено Л. П. $L[g]$ - целая функция экспоненциального типа. Подробнее, пусть



\[

E(\omega)=\underset{K \subset \omega}{\operatorname{limind} \lim } \operatorname{proj}_{\varepsilon} E_{\varepsilon}(K),

\]



где $K$ - вещественный компакт, $E_{\varepsilon}(K)$ - банахово пространство всех целых функций $f(z)$ с конечной нормой



\[

|f|=\sup \left\{|f(z)| e^{-H_{K}(z)-\varepsilon|z|} ; z \in \mathbb{C}^{n}\right\}

\]



где



\[

H_{K}(z)=\sup \{-y \zeta ; \zeta \in K\}, z=x+i y,

\]



поскольку компакт $K$ вещественный. Тогда $L$ - линейное непрерывное инъективное отображение .d्' $(\omega)$ в $E(\omega)$, а в случае выпуклой области $\omega, L-$ топологич. изоморфизм.



Л. П. обобщенных функций с компактным н о с и те л е м. Основное пространство $\mathscr{E}(\omega), \omega-$ область в $\mathbb{R}^{n}$, также содержит все линейные экспоненты, так что для всякой $g \in \mathscr{E}^{\prime}(\omega)$ определено Л. П. $L[g]-$ целая функция экспоненциального типа. Подробнее, пусть



\[

E_{\infty}(\omega)=\lim \text { ind } E_{p}(K),

\]



где $E_{p}(K)$ - банахово пространство всех целых функций $f(z)$ с конечной нормой



\[

|f|=\sup \left\{|f(z)| e^{-H_{K}(z)}(1+|z|)^{-p} ; z \in \mathbb{C}^{n}\right\},

\]



индуктивный предел берется по всем $p=0,1,2, \ldots$, и компактам $K \subset \omega$. Тогда $L-$ линейное непрерывное инъек-





\end{document}